\chapter{Introducción}
\label{sec.tp2:introduccion}

La iluminación en los ámbitos laborales constituye uno de los factores ambientales de mayor relevancia para la
seguridad, salud y productividad de los trabajadores. Un sistema de iluminación inadecuado puede generar fatiga visual,
cefaleas, errores operativos e incluso incrementar significativamente la tasa de siniestralidad. Por el contrario, un
nivel y calidad óptimos de iluminación favorecen el desempeño cognitivo y motor, disminuyendo el cansancio y previniendo
accidentes.

En la República Argentina, el marco normativo establece criterios técnicos estrictos para asegurar condiciones dignas
de iluminación. El Decreto N° 351/79, reglamentario de la Ley de Higiene y Seguridad en el Trabajo N° 19.587,
determina las intensidades mínimas de iluminación requeridas según el tipo de edificio, local y tarea visual en su
Anexo IV. Asimismo, la Resolución SRT N° 84/2012 estipula un protocolo de carácter obligatorio para registrar y evaluar
las mediciones de iluminación en los puestos de trabajo.

\section{Objetivos}
\label{sec.tp2:objetivos}

Los objetivos principales de este trabajo práctico son:
\begin{itemize}
    \item Relevar el sistema de iluminación actual en un ambiente real seleccionado.
    \item Calcular las dimensiones del local y determinar la grilla de medición conforme a las exigencias normativas de la Resolución SRT N° 84/2012.
    \item Medir los niveles de iluminancia utilizando un instrumento de medición adecuado (luxómetro o aplicación calibrada).
    \item Determinar la iluminancia media ($E_{media}$) y verificar la relación de uniformidad del local.
    \item Comparar los resultados con las exigencias del Decreto N° 351/79 para evaluar la conformidad.
    \item Elaborar conclusiones técnicas y, si corresponde, proponer mejoras en la instalación.
\end{itemize}
